% Interfaz computable CG-008; formalización nueva de un lector residual.
% Productor: los caracteres e intervalos regionales construidos previamente.
% No identifica la componente numérica con el estado enriquecido completo.

\section{Computable residual initialization and compatibility of regional emissions}
\label{sec:residuo-regional-computable}

The numerical realization of a regional section admits an executable
description by rational intervals. This description allows the residue
to be initialized, updated and its blocks recovered without introducing a previously
fixed decimal string. Its origin is the regional producer of
\cref{x_reciprocidad:sec:generacion}: the character,
its orientation and its normalization are first determined; its rational
approximations are then constructed. The present development formalizes the residual component of that
reading. The trajectory, sheet, incidences and memory of the enriched
state simultaneously preserve their own operations.

The order of operations is therefore
\[
 \begin{gathered}
 \text{generated regional state}
 \longrightarrow \text{rational intervals of the character}
 \\
 \longrightarrow \text{rational name of the residue}
 \longrightarrow \text{compatible positional emissions}.
 \end{gathered}
\]
Comparison with the Archimedean realization appears in the correctness
proof. The algorithm receives intervals; none of its instructions
receives a conventional real number or a list of target digits.

\subsection{Rational name produced by a regional reader}
\label{subsec:residual-nombre-racional}

A region is fixed and \(j\mapsto[\ell_j,u_j]\) denotes its rational
producer. The interface conditions are
\begin{equation}
 \ell_j\leq v\leq u_j,\qquad
 \lim_{j\to\infty}(u_j-\ell_j)=0,
 \qquad \ell_j,u_j\in\mathbb Q.
 \label{eq:residual-contrato-productor}
\end{equation}
Here \(v\) designates the exact realization of the character to formulate
correctness; the procedure calculating \(\ell_j,u_j\) comes from the regional
reader. The closure, propagation and autoscale intervals already
constructed satisfy these conditions. In the implementation, its interface
is a function that, given \(j\), returns two exact fractions.

Nesting is obtained, if necessary, through the finite operation
\begin{equation}
 L_j=\max_{0\leq i\leq j}\ell_i,
 \qquad U_j=\min_{0\leq i\leq j}u_i,
 \qquad I_j=[L_j,U_j].
 \label{eq:residual-anidacion}
\end{equation}

\begin{proposition}[Nesting and preservation of the realization]
\label{prop:residual-anidacion}
The intervals \(I_j\) are rational, nonempty and nested. They contain
\(v\), their width tends to zero and \(\bigcap_j I_j=\{v\}\).
\end{proposition}
\begin{proof}
Each lower bound is less than or equal to \(v\), and each upper bound is
greater than or equal to \(v\); the same property holds for the finite maximum and
minimum. As \(j\) increases, the maximum increases and the minimum decreases.
Moreover, \(0\leq U_j-L_j\leq u_j-\ell_j\). If two points belonged to
the intersection, their distance would be bounded by each width and would be zero.
Membership of \(v\) establishes nonemptiness.
\end{proof}

\subsection{Initial prefix and representation of the residue}
\label{subsec:residual-inicializacion}

Let \(b\geq2\) be an integer base and \(n_0\geq0\) a depth.
The realization is assumed to avoid the corresponding boundaries:
\begin{equation}
 b^n v\notin\mathbb Z\qquad(n\geq0).
 \label{eq:residual-no-frontera}
\end{equation}
The irrationality of the three regional values implies this property;
its proof or use as a hypothesis belongs to the fixed reader.
For a concrete depth, avoidance at the scales used
suffices. Realizations with a terminating expansion additionally require
an exact boundary decider and a representation convention.

We define the rational predicate
\[
 \mathsf{Sep}(j,n):\quad
 \lfloor b^n L_j\rfloor=\lceil b^n U_j\rceil-1.
\]
It is decidable through integer operations on numerators and
denominators. Let
\begin{equation}
 j_0=\min\{j:\mathsf{Sep}(j,n_0)\},
 \qquad N_0=\lfloor b^{n_0}L_{j_0}\rfloor.
 \label{eq:residual-prefijo-inicial}
\end{equation}
The integer \(N_0\) is calculated from the producer. Its use as an antecedent
of the residual representation does not require introducing it as a literal datum.

\begin{theorem}[Computable initialization of the regional residue]
\label{thm:residual-inicializacion}
The index \(j_0\) exists. For \(j\geq j_0\), the intervals
\begin{equation}
 E_{0,j}=
 \bigl[b^{n_0}L_j-N_0,\ b^{n_0}U_j-N_0\bigr]
 \label{eq:residual-nombre-inicial}
\end{equation}
are nonempty, nested and contained in \([0,1]\). Their width tends to
zero and their intersection consists of a single element
\(\varepsilon_0\in(0,1)\). In the Archimedean realization one has
\begin{equation}
 N_0=\lfloor b^{n_0}v\rfloor,
 \qquad b^{n_0}v=N_0+\varepsilon_0.
 \label{eq:residual-identificacion-inicial}
\end{equation}
The rational sequence \((E_{0,j})_{j\geq j_0}\) constitutes a
computable representation of the residue.
\end{theorem}
\begin{proof}
The point \(b^{n_0}v\) lies in the interior of a cell with integer
endpoints. Its distance to those endpoints is positive. When the width of
\(b^{n_0}I_j\) is less than that distance, both endpoints belong to the
same cell and \(\mathsf{Sep}(j,n_0)\) holds. There is therefore a
first separation index. At any separated index, the integer
\(\lfloor b^{n_0}v\rfloor\) lies between the two equal integers
of the predicate and coincides with them. The initial equality places the endpoints
of \(b^{n_0}I_{j_0}\) between \(N_0\) and \(N_0+1\). Nesting
preserves that inclusion. The affine transformation of
\eqref{eq:residual-nombre-inicial} preserves nesting and nonemptiness and
multiplies the widths by \(b^{n_0}\). Their intersection is the point of
\eqref{eq:residual-identificacion-inicial}; boundary avoidance
places it in \((0,1)\). Each endpoint is calculated through exact fractions,
and the search terminates by the preceding argument.
\end{proof}

The identity \eqref{eq:residual-identificacion-inicial} is a correctness
conclusion. The executable representation is the sequence
\eqref{eq:residual-nombre-inicial}, constructed from regional
intervals. This distinction separates the production of a rational name from the
recognition of its value.

\subsection{Residual update and emission through exact arithmetic}
\label{subsec:residual-actualizacion}

Suppose the reading state \((n_0+k,N_k,j_k)\) has been constructed, together
with the name
\begin{equation}
 E_{k,j}=
 \bigl[b^{n_0+k}L_j-N_k,\ b^{n_0+k}U_j-N_k\bigr],
 \qquad j\geq j_k.
 \label{eq:residual-nombre-k}
\end{equation}
One increases \(j\) from \(j_k\) until
\[
 \left\lfloor b\inf E_{k,j}\right\rfloor
 =\left\lceil b\sup E_{k,j}\right\rceil-1.
\]
The first index satisfying this equality is denoted by \(j_{k+1}\), and the
common integer is denoted by \(d_k\). The update is
\begin{equation}
 N_{k+1}=bN_k+d_k,
 \qquad E_{k+1,j}=bE_{k,j}-d_k
 \quad(j\geq j_{k+1}).
 \label{eq:residual-actualizacion-exacta}
\end{equation}
Multiplication of an interval by \(b>0\) and subtraction of \(d_k\)
act on its two endpoints. The reading index reached,
the depth and the emitted block are preserved.

\Needspace{20\baselineskip}
\begin{theorem}[Termination at arbitrary depth and positional commutation]
\label{thm:residual-compatibilidad-universal}
Under \eqref{eq:residual-contrato-productor} and
\eqref{eq:residual-no-frontera}, each update terminates and, for every
\(k\geq0\), one has
\begin{align}
 N_k&=\lfloor b^{n_0+k}v\rfloor,
 &0&\leq d_k<b,\label{eq:residual-prefijos-k}\\
 \varepsilon_{k+1}&=b\varepsilon_k-d_k,
 &d_k&=\lfloor b\varepsilon_k\rfloor,\label{eq:residual-recurrencia}\\
 b^{n_0+k}v&=N_k+\varepsilon_k,
 &0&<\varepsilon_k<1.\label{eq:residual-invariante}
\end{align}
If \(P(b,n)\) is the emission by rational separation of
\cref{thm:df-emision-correcta}, then
\begin{equation}
 N_k=P(b,n_0+k),\qquad
 d_k=P(b,n_0+k+1)\bmod b.
 \label{eq:residual-igualdad-publicacion}
\end{equation}
The stopping indices of the two implementations may differ;
their emitted integers coincide exactly.
\end{theorem}
\begin{proof}
Initialization provides the assertions for \(k=0\). Assuming them
at level \(k\), multiplying the invariant by \(b\) gives
\[
 b^{n_0+k+1}v=bN_k+b\varepsilon_k.
\]
Boundary avoidance implies \(b\varepsilon_k\notin\mathbb Z\).
The intervals \(bE_{k,j}\) are nested, contain that point and have
width tending to zero. The separation argument of the preceding theorem
guarantees termination and fixes
\(d_k=\lfloor b\varepsilon_k\rfloor\). Since
\(0<\varepsilon_k<1\), the block belongs to
\(\{0,\ldots,b-1\}\). Subtracting \(d_k\) transforms the intervals into
a nested name of the next residue and preserves its membership in
\((0,1)\). The integer part of the preceding equality is
\(bN_k+d_k\), which proves the new invariant. Induction holds for
every finite depth, without imposing a maximum depth.
Finally, the rational emission theorem also identifies
\(P(b,n)\) with \(\lfloor b^n v\rfloor\). Equality of the
prefixes and the residue formula modulo \(b\) prove
\eqref{eq:residual-igualdad-publicacion}.
\end{proof}

\begin{corollary}[Truncation, recovery and residual bounds]
\label{cor:residual-recuperacion}
For \(k,h\geq0\),
\[
 \left\lfloor\frac{N_{k+h}}{b^h}\right\rfloor=N_k,
 \qquad N_k=\frac{N_{k+1}-d_k}{b},
 \qquad \varepsilon_k=\frac{\varepsilon_{k+1}+d_k}{b}.
\]
The width of the residual name at index \(j\) is
\(b^{n_0+k}(U_j-L_j)\). Consequently, any positive rational precision
for the residue is reached after a finite number of queries to the
producer. The depth and the archived blocks determine recovery
in the reading component.
\end{corollary}
\begin{proof}
Division of the extended prefix recovers the preceding prefix by
\eqref{eq:residual-prefijos-k}. The two recovery formulas are
obtained by rearranging \eqref{eq:residual-actualizacion-exacta} and
\eqref{eq:residual-recurrencia}. The width formula follows immediately from
\eqref{eq:residual-nombre-k}; its convergence to zero provides
termination of the precision search.
\end{proof}

\subsection{Correspondence between ternary and decimal depths}
\label{subsec:residual-bases}

In fractional normalization one writes \(v=m+r\), where
\(m=P(b,0)\) is obtained from the same producer. The fractional prefix is
\(C_n=P(b,n)-mb^n\). By integer calculation,
\[
 b^n v-P(b,n)=b^n r-C_n.
\]
This equality reconciles the prefix including the integer part with the one
starting at zero, preserving an identical residue. Subtracting \(m\)
is a subsequent transformation of the regional name.

Tritic grouping uses \(b=729=3^6\). Six blocks correspond to
\(36\) trits; the specialization \(n_0=6\) provides a rational
name of the residue at that depth. Decimal grouping uses
\(b=1000=10^3\); six blocks correspond to \(18\) decimal digits.
They are different reading coordinates of the same character. Their compatibility
is expressed by the cylinders containing that realization and by the
scale-change identities; it does not require equating their numerical residues.

The residual antecedent of
\cref{sec:c27-generacion-coinductiva-profundidad-arbitraria} distinguishes the
internal state from its publication. The present interface provides an
executable rational representation of the numerical component, with the same
residual recurrence once its reading is fixed. To identify it with another
residual realization of that section, it suffices to establish that both read the same
character, at the same depth and with the same normalization: the identity
\(b^n v=N_n+\varepsilon_n\) and uniqueness of the prefix then force
equality of the residues. This comparison is an equality of readers.
The evolution of positions, sheets, incidence and memory remains governed
by enriched TPK and by their respective transport theorems.

\subsection{Exact implementation contract}
\label{subsec:residual-implementacion}

The implementation of this interface uses exclusively rational
fractions, integer powers, maximum, minimum, floor division and ceiling
division. The interval producer is an explicit argument of the
program. A query does not return \(v\): it returns its two generated bounds.
The execution state preserves base, depth, prefix, index reached
and ordered emission register. For each step, the register stores the
preceding residual name at the decisive index, the block and the transformed
name; the two recoveries in the corollary are checked through
rational equalities.

The preceding universal proofs justify termination and
compatibility under the producer contract. Finite tests verify
the implementation of those operations and their identities; the generality
of the result lies in induction and in the rational separation of cells.
This formalization combines a computable reading and recovery interface
without replacing regional selection or attributing to the numerical component
the additional operations of the enriched state.
